% !TeX root = ../navier-stokes-manuscript.tex
\subsection{Critical height and the Zeno cascade}\label{sub:CriticalHeightAndTheZenoCascade}

Rescale the vortex before deciding what its narrowing means. For \(\lambda>0\), the Navier--Stokes scaling is
\begin{equation}\label{eq:BodyNavierStokesScaling}
  u_\lambda(x,t)=\lambda u(\lambda x,\lambda^2t),
  \qquad
  p_\lambda(x,t)=\lambda^2p(\lambda x,\lambda^2t).
\end{equation}
Every term acquires the same factor:
\begin{equation}\label{eq:BodyBalancedEquationScaling}
\begin{aligned}
  \partial_tu_\lambda
    &=\lambda^3(\partial_tu)(\lambda x,\lambda^2t),
  &
  (u_\lambda\cdot\nabla)u_\lambda
    &=\lambda^3((u\cdot\nabla)u)(\lambda x,\lambda^2t),
  \\
  \nabla p_\lambda
    &=\lambda^3(\nabla p)(\lambda x,\lambda^2t),
  &
  \nu\Delta u_\lambda
    &=\lambda^3(\nu\Delta u)(\lambda x,\lambda^2t).
\end{aligned}
\end{equation}
At a smaller scale, viscosity has not gained a relative advantage over nonlinear deformation, and nonlinear deformation has not gained one over viscosity. The same balance repeats. This is what makes the perfectly stretching vortex a critical picture rather than a competition settled by scaling alone.

Now compare that balance with kinetic energy. Pair \eqref{eq:NavierStokesSystem} with \(u\). Incompressibility removes transport and pressure, while integration by parts turns viscosity into dissipation:
\begin{equation}\label{eq:BodyEnergyIdentity}
  \frac12\norm{u(t)}_2^2
  +\nu\int_0^t\norm{\nabla u(s)}_2^2\,ds
  =\frac12\norm{u_0}_2^2.
\end{equation}
The \(L^2\) energy stays finite, and one spatial derivative is square-integrable after we integrate in time. The identity gives no pointwise-in-time ceiling for the gradient. Under the same scaling, the relevant norms move in opposite directions:
\begin{equation}\label{eq:BodyScalingNormComparison}
  \begin{aligned}
    \norm{u_\lambda(t)}_2
      &=\lambda^{-1/2}\norm{u(\lambda^2t)}_2,\\
    \norm{u_\lambda(t)}_\infty
      &=\lambda\norm{u(\lambda^2t)}_\infty,\\
    \norm{\nabla u_\lambda(t)}_2
      &=\lambda^{1/2}\norm{\nabla u(\lambda^2t)}_2.
  \end{aligned}
\end{equation}
A finer copy carries less kinetic energy while its amplitude and gradient grow. The vortex can sharpen without violating the energy ceiling in \eqref{eq:BodyEnergyIdentity}.

For homogeneous Sobolev norms,
\[
  \norm{u_\lambda(t)}_{\dot H^s}
  =\lambda^{s-1/2}\norm{u(\lambda^2t)}_{\dot H^s}.
\]
The exponent vanishes at \(s=1/2\). Write \(\Lambda:=(-\Delta)^{1/2}\) and measure the solution at that height by
\begin{equation}\label{eq:BodyCriticalHeightDefinition}
  H(t)
  :=\frac12\norm{u(t)}_{\dot H^{1/2}}^2
  =\frac12\norm{\Lambda^{1/2}u(t)}_2^2.
\end{equation}
The embedding \(\dot H^{1/2}(\R^3)\hookrightarrow L^3(\R^3)\) now transfers the endpoint escape into the same scale-invariant quantity:
\begin{equation}\label{eq:CriticalHeightMustEscape}
  T_*<\infty
  \quad\Longrightarrow\quad
  \sup_{0<t<T_*}H(t)=\infty.
\end{equation}

A fixed profile can be pushed to finer frequencies by scaling without changing \(H\) at all. The actual finite-time history cannot stop there: \eqref{eq:CriticalHeightMustEscape} forces its critical height past every finite level. This is the distinction we needed. Scale tells us how narrow the vortex is; critical height tells us whether the complete solution has acquired more scale-invariant content.

Once \(H\) has to pass every finite level before \(T_*\), its time geometry is forced. The solution is classical before the endpoint, so \(H\) is continuous on \([0,T]\) for every \(T<T_*\). Choose \(H_0>H(0)\), set \(L_k=2^kH_0\), and let
\[
  t_0:=\inf\{t\in(0,T_*):H(t)=H_0\},
\]
and, once \(t_k\) is fixed, let
\[
  t_{k+1}:=\inf\{t\in(t_k,T_*):H(t)=L_{k+1}\}.
\]
Every set is nonempty by \eqref{eq:CriticalHeightMustEscape}, and continuity places the stated level at its first-passage time. Hence
\begin{equation}\label{eq:ZenoFirstPassageTimes}
  t_0<t_1<t_2<\cdots<T_*,
  \qquad
  H(t_k)=2^kH_0,
  \qquad
  t_k\uparrow T_*.
\end{equation}
The times cannot accumulate earlier: any earlier limit would lie inside the smooth interval, where \(H\) is locally bounded. They can accumulate only at \(T_*\).

This is the Zeno cascade. The same fluid crosses infinitely many doubled critical-height levels before one finite last time. Nothing says that \(H\) rises monotonically between those crossings; the construction records only the first time each new level is reached. The velocity has now been forced through infinitely many critical events in finite time. Repeated integration tells us what that demands from its temporal responses.
